% Compiles on its own (pulls ../main.tex's preamble via subfiles) or as part of the book.
\documentclass[../main.tex]{subfiles}
\begin{document}

\chapter{Augmentation Without Breaking the Pair}

This chapter is the theory gate for the last piece of the data pipeline: data augmentation,
and the unit test that keeps it honest. It is a short chapter about a large class of silent
bugs.

\section{Why augmentation is dangerous here specifically}

Data augmentation --- randomly flipping, cropping, scaling, and colour-jittering training
images --- is standard and usually harmless. In single-image detection you transform the image
and its boxes together, and if you get it slightly wrong you notice, because the boxes visibly
drift off the objects.

Multispectral fusion removes that safety net. We now transform \emph{two} images that must stay
in pixel-for-pixel correspondence, because middle fusion adds their feature maps together at a
matching spatial location. If a random crop takes a slightly different window from the RGB image
than from the thermal image, then at fusion time we are adding the features of a pedestrian in
one stream to the features of the pavement beside them in the other. The network still trains.
The loss still goes down. The model is just quietly worse, and nothing tells you why.

\begin{lesson}
The correspondence that fusion depends on is created by the loader and can be destroyed by the
augmentation. There is no error message for a broken pair --- only a worse number that looks like
a bad architecture. So we do not \emph{trust} the augmentation to preserve correspondence; we
\emph{test} it, once, with an assertion that fails loudly the day someone breaks it.
\end{lesson}

\section{The rule, and why it has two halves}

Augmentations split cleanly into two kinds, and the two kinds are treated oppositely.

\paragraph{Geometric augmentations move pixels.} Horizontal flip, crop, scale, rotation,
translation. These \emph{must} be applied identically --- same operation, same random parameters
--- to the RGB image, the thermal image, \emph{and} the boxes. The whole point is that all three
move together, so the correspondence survives. Sample the random flip decision \emph{once} and
feed it to all three; never let each stream roll its own dice.

\paragraph{Photometric augmentations change values, not positions.} Brightness, contrast, hue,
saturation jitter. These are applied to the \emph{RGB stream only}. Two reasons. First, they do
not affect correspondence (a pixel keeps its location), so there is no pairing constraint.
Second, and more importantly, they are \emph{physically meaningless on thermal}. A thermal pixel
is a radiometric measurement --- roughly, an apparent temperature. Randomly scaling its
``brightness'' or shifting its ``hue'' corrupts the one quantity the thermal stream exists to
carry. Colour-jittering a thermal image is not augmentation; it is noise injection into your most
reliable night-time signal.

\begin{lesson}
Geometric augmentations: identical on both streams, boxes included, one shared random draw.
Photometric augmentations: RGB only. Getting the split backwards --- jittering thermal, or
letting the two streams flip independently --- is among the most common and most invisible causes
of ``my fusion model is worse than my single-modality baseline.''
\end{lesson}

\section{Testing correspondence, not trusting it}

We prove the geometric path with a constructed example rather than an inspected one. Place an
identical bright marker at a known pixel in \emph{both} synthetic modalities, and a box around
it. Apply a paired geometric transform. Then assert two things:

\begin{enumerate}
  \item the marker lands at the \emph{same} pixel location in the transformed RGB and the
        transformed thermal --- the two streams moved together; and
  \item the transformed box still contains the transformed marker --- the boxes moved with the
        pixels.
\end{enumerate}

If someone later edits the transform so that only one stream flips, or the boxes are left
untouched, the first or second assertion fails immediately, at test time, instead of six weeks
later as an unexplained dip in the results table. This is the first \emph{unit test} of the
project: a small, permanent, automatic check of a property we care about, run by \texttt{pytest}.
It costs a minute to write and it retires an entire category of silent failure.

\section{Scope for now}

We do not need the full training-time augmentation zoo yet (Mosaic, MixUp, and the rest arrive
with the training phase, and each will extend the same correspondence test). For now we implement
the two representatives that exercise both halves of the rule --- a paired horizontal flip
(geometric) and an RGB-only brightness/contrast jitter (photometric) --- and the test that guards
them. The lesson is the invariant, not the size of the menu.

\end{document}
